% s-164-p164-m001; original PDF 164, 165
\Aphorism{s-164-p164-m001}{250}
Was ist Jugend für ein erbärmlicher Zustand! Dieser „Drangdruck“ der sich entwickelnden und einander hemmenden Organe, dieses Durcheinander ohne Polyphonie, diese Schmerzen des Wachsthums, diese unbehaglichen Zerrungen und Verrenkungen der Seele, diese geistige Mutation, bei der man nicht Herr der eigenen Stimme ist und lauter falsche Töne herausbringt! Der Anspruch, für voll genommen zu werden, zusammen mit dem bösen Gewissen, noch lange nicht voll zu sein, die ausgehängte Überfertigkeit des Urtheils bei inwendigem Schwanken und Nicht-aus-noch-ein-wissen, die grosse Gebärde und Leidenschaft ohne den grossen Stil! Barbarische Unempfindlichkeit einerseits (was muthet solch ein junger Idealist, gesetzt dass er deutscher Student ist, allein seinem \Sperr{Magen} zu!) mit reizbarer Hyperästhesis andererseits: wie quält so Einen der Anblick des Wolkenhimmels oder eines Mädchengesichts — ein wahres Sturmläuten in den Nerven, ohne dass die intellectuelle Glocke zu tönen käme! Geschmacklose Eklexis: alles wird versucht, alles wieder verworfen, der kleine Geist flattert um die grössten Probleme, jeder Tag bringt einen neuen —ismus, bis es zu einem Damaskus kommt und der Apostel seinen Heiland findet. Nun ebenso geschmacklose Verehrung und Dogmatik, mit wüthender Selbstblendung gegen alles „Andere“ und fatalem Bekehrungseifer an „Anderen“; ein \Foreign{credo quia absurdum}, zu deutsch Verliebtheit, die sich mit besonderer Inbrunst auf die Schwächen des geliebten Gegenstandes wirft. Es thut einem kühleren und genussfähigeren Geiste wirklich leid, solch einen verliebten Wirr- und Brausekopf sich in eine sterile Monomanie verbeissen zu sehen, wo doch der geniessbaren Dinge so viele sind; aber er wage nicht, ihm den Knochen aus den Zähnen zu reissen und Fleisch dafür zu geben, wenn er die Dogge Enthusiasmus nicht am Halse haben will. — Nein, Jugend mag Tugend haben, soviel sie behauptet, aber Geschmack, Stil, Vernunft hat sie nicht. Jugend ist ein Martyrium, eine Bartholomäusschindung, ein Laurentiusrost, jedenfalls nichts Wohlriechendes.
\par

